\setcounter{chapter}{-1}
\appendiceschapter{รายการ API หลักของระบบ}

ภาคผนวกนี้สรุป API ที่ใช้เชื่อมส่วนหน้ากับส่วนหลังในซอร์สโค้ดปัจจุบัน เพื่อใช้เป็นแนวทางสำหรับการทดสอบแบบ end-to-end และการพัฒนาต่อไป

\begin{table}[H]
\centering
\caption{รายการ API หลักของ Blur System}
\label{tab:api-list}
\begin{tabular}{|p{0.16\textwidth}|p{0.39\textwidth}|p{0.30\textwidth}|}
\hline
\textbf{Method} & \textbf{Endpoint} & \textbf{หน้าที่} \\
\hline
POST & \texttt{/api/upload} & รับภาพหรือวิดีโอและสร้างข้อมูลเมตา \\
\hline
GET & \texttt{/api/preview-frames/\{video\_id\}} & อ่านเฟรมตัวอย่างสำหรับเลือกเฟรม \\
\hline
POST & \texttt{/api/start-detection/\{video\_id\}} & เริ่มงานตรวจจับใบหน้าแบบเบื้องหลัง \\
\hline
GET & \texttt{/api/status/\{video\_id\}} & อ่านสถานะและความคืบหน้าของงานตรวจจับ \\
\hline
GET & \texttt{/api/recognized-persons/\{video\_id\}} & อ่านกลุ่มบุคคลที่ตรวจพบสำหรับ Face Map \\
\hline
POST & \texttt{/api/update-blur-selection} & บันทึกการเลือกปกปิดรายบุคคล \\
\hline
GET/POST & \texttt{/api/manual-blur/\{video\_id\}} & อ่านหรือบันทึกกรอบ Manual Blur \\
\hline
POST & \texttt{/api/render-preview/\{video\_id\}} & สร้างตัวอย่างผลลัพธ์ \\
\hline
POST & \texttt{/api/export-video/\{video\_id\}} & สร้างและส่งออกไฟล์ผลลัพธ์ \\
\hline
\end{tabular}
\end{table}

\section*{ขั้นตอนใช้งานโดยสังเขป}

\begin{enumerate}[label=(\arabic*), leftmargin=*]
    \item อัปโหลดภาพหรือวิดีโอที่ระบบรองรับ
    \item เลือกโหมด Auto Blur, Detect \& Select หรือ Manual Control
    \item หากใช้ Detect \& Select ให้รอการตรวจจับ แล้วเลือกบุคคลที่ต้องการปกปิดหรือยกเว้น
    \item ปรับชนิดและความเข้มของการปกปิด หรือวาดกรอบ Manual Blur เพิ่มเติม
    \item สร้างตัวอย่างเพื่อตรวจสอบความครอบคลุมของการปกปิด แล้วจึงส่งออกไฟล์ผลลัพธ์
\end{enumerate}
